% !TEX program = platex
% 数学ノート集・共通テンプレート v2（積分論2026年度第1部の版面を参考）
% A5判。章内で定義・定理・例・注意等の番号を共通化。
% コンパイル： platex -interaction=nonstopmode math_note_template_v2.tex (2回)
%              dvipdfmx math_note_template_v2.dvi
\documentclass[a5paper,10pt,onecolumn]{jsarticle}
\usepackage[dvipdfmx,top=22mm,bottom=20mm,inner=19mm,outer=18mm,headheight=13pt,headsep=6mm]{geometry}
\usepackage{amsmath,amssymb,amsthm,mathtools}
\usepackage{enumitem,fancyhdr,titlesec}
\usepackage[dvipdfmx,bookmarks=true,colorlinks=false]{hyperref}
\usepackage{pxjahyper}

% ----- 差し替える文書情報 -----
\newcommand{\NoteTitle}{格子・非格子の局所中心極限定理}
\newcommand{\NoteSeries}{数学ノート集}

% ----- 文字・組版 -----
% 参考PDFの質素な明朝＋Computer Modernの印象を引き継ぎながら，
% A5判で窮屈にならないよう本文の行間を少し確保する．
\linespread{1.13}
\setlength{\parindent}{1zw}
\setlength{\parskip}{0pt}
\setlength{\abovedisplayskip}{6pt plus 2pt minus 2pt}
\setlength{\belowdisplayskip}{6pt plus 2pt minus 2pt}
\setlength{\abovedisplayshortskip}{3pt plus 2pt minus 1pt}
\setlength{\belowdisplayshortskip}{4pt plus 2pt minus 1pt}
\setlist[enumerate]{topsep=3pt,itemsep=1pt,parsep=0pt,leftmargin=2.4zw}
\titleformat{\section}{\normalfont\normalsize\bfseries}{\thesection}{0.8zw}{}
\titlespacing*{\section}{0pt}{2.0ex plus .4ex}{.9ex}
\titleformat{\subsection}{\normalfont\normalsize\bfseries}{\thesubsection}{0.7zw}{}
\titlespacing*{\subsection}{0pt}{1.5ex plus .3ex}{.6ex}
\raggedbottom

% ----- 柱とページ番号：ページ上部、罫線なし -----
\fancypagestyle{mathnote}{%
  \fancyhf{}%
  \fancyhead[L]{\footnotesize \NoteSeries\ （\NoteTitle）}%
  \fancyhead[R]{\footnotesize \thepage}%
  \renewcommand{\headrulewidth}{0pt}%
  \renewcommand{\footrulewidth}{0pt}%
}
\pagestyle{mathnote}

% ----- 節内通し番号の定義・定理群 -----
\newtheoremstyle{mathnotetheorem}
 {4.5pt}{4.5pt}{\normalfont}{}%
 {\bfseries}{}{0.65zw}%
 {\thmname{#1}\thmnumber{ #2}\thmnote{（#3）}}
\theoremstyle{mathnotetheorem}
\newtheorem{theorem}{定理}[section]
\newtheorem{definition}[theorem]{定義}
\newtheorem{lemma}[theorem]{補題}
\newtheorem{proposition}[theorem]{命題}
\newtheorem{corollary}[theorem]{系}
\newtheorem{example}[theorem]{例}
\newtheorem{remark}[theorem]{注意}
\renewcommand{\proofname}{証明}
\renewcommand{\qedsymbol}{$\square$}
\makeatletter
\renewenvironment{proof}[1][\proofname]{%
  \par\pushQED{\qed}\normalfont\topsep4\p@\@plus2\p@\relax
  \trivlist\item[\hskip\labelsep\bfseries #1]\ignorespaces
}{\popQED\endtrivlist\@endpefalse}
\makeatother

% ----- 使うノートで適宜拡張する記号 -----
\newcommand{\R}{\mathbb{R}}
\newcommand{\norm}[1]{\lVert #1\rVert}
\newcommand{\abs}[1]{\lvert #1\rvert}

\begin{document}
\section*{読了後にできること}
本稿を読むための前提は，二項分布，Stirlingの公式，局所中心極限定理．
本稿の具体的な到達課題は，格子確率の一点確率をStirlingの公式から評価する．


\par\noindent\textbf{記号と約束.}\quad
関数空間の添字 $c$ はコンパクトな台をもつことを表す．
\par

\section{問題設定と結論}
確率変数 $X_1,X_2,\ldots$ を独立同分布とし，
$S_n=X_1+\cdots+X_n$ とする．通常の中心極限定理は
$S_n$ の分布函数の極限を述べるが，本稿では
各点の確率や密度の収束を扱う．後者には分布の算術的な性質が必要である．

\begin{definition}
$\mathbb Z^d$ 値確率変数 $X$ を\emph{強非周期的}というのは，
$\{a-b:a,b\in\operatorname{supp}(X)\}$ の生成する
加法部分群が $\mathbb Z^d$ 全体である場合をいう．
ここで $\operatorname{supp}(X)=\{a\in\mathbb Z^d:\Pr(X=a)>0\}$ とする．
\end{definition}

\begin{theorem}[多次元格子局所中心極限定理]\label{thm:llclt}
$X_1,X_2,\ldots$ を独立同分布な $\mathbb Z^d$ 値確率変数とし，
$\mathbb E|X_1|^2<\infty$ とする．$X_1$ は強非周期的であり，
平均 $\mu=\mathbb E X_1\in\mathbb R^d$，共分散行列
$\Sigma=\mathbb E[(X_1-\mu)(X_1-\mu)^T]$ が正定値であると仮定する．
$z_{n,k}=\Sigma^{-1/2}(k-n\mu)/\sqrt n$ とおくと
\[
 \sup_{k\in\mathbb Z^d}\left|
 n^{d/2}\sqrt{\det\Sigma}\,\Pr(S_n=k)
 -(2\pi)^{-d/2}e^{-|z_{n,k}|^2/2}
 \right|\longrightarrow0.
\]
\end{theorem}

ここでの収束は $k$ に関して一様である．特定の点 $k=k_n$ が
$n\mu+O(\sqrt n)$ にある場合だけでなく，すべての格子点を同時に扱う．

\section{Fourier反転と算術的非周期性}
\begin{lemma}[格子上のFourier反転]\label{lem:inv}
$Y$ を $\mathbb Z^d$ 値確率変数とし，
$\varphi_Y(t)=\mathbb E(e^{it\cdot Y})$ とおく．
任意の $k\in\mathbb Z^d$ に対して
\[
 \Pr(Y=k)=(2\pi)^{-d}\int_{[-\pi,\pi]^d}
             e^{-it\cdot k}\varphi_Y(t)\,dt
\]
である．
\end{lemma}
\begin{proof}
各 $x,k\in\mathbb Z^d$ について
\[
 (2\pi)^{-d}\int_{[-\pi,\pi]^d}e^{it\cdot(x-k)}\,dt
 =\prod_{j=1}^d\frac1{2\pi}\int_{-\pi}^{\pi}
         e^{it_j(x_j-k_j)}\,dt_j
 =\mathbf1_{\{k\}}(x)
\]
となる．積分内の絶対値は1だから，Fubiniの定理によって
期待値と積分を交換し，$x=Y$ として結論を得る．
\end{proof}

\begin{lemma}[単位円上で絶対値が1となる条件]\label{lem:unit}
$Y$ を任意の確率変数，$Z=e^{it\cdot Y}$ とする．
$|\mathbb EZ|=1$ であることと，$Z$ が確率1である
定数に一致することは同値である．
\end{lemma}
\begin{proof}
$|Z|=1$ である．$m=\mathbb EZ$ とすれば
\[
 \mathbb E|Z-m|^2=\mathbb E|Z|^2-|m|^2=1-|m|^2.
\]
ゆえに $|m|=1$ のとき右辺は0となり，$Z=m$ が確率1で成立する．
逆向きは明らかである．
\end{proof}

\begin{lemma}[強非周期性のFourier的特徴付け]\label{lem:period}
$\mathbb Z^d$ 値確率変数 $X$ が強非周期的であることと
\[
 |\varphi_X(t)|<1\qquad
  (t\in[-\pi,\pi]^d\setminus\{0\})
\]
が同値である．
\end{lemma}
\begin{proof}
$|\varphi_X(t)|=1$ なら前補題により
$e^{it\cdot a}=e^{it\cdot b}$ が
すべての $a,b\in\operatorname{supp}(X)$ について成立する．
つまり $t\cdot(a-b)\in2\pi\mathbb Z$ である．
差の生成する群が $\mathbb Z^d$ なら
$t_j\in2\pi\mathbb Z$ が全ての $j$ に対して成立する．
$[-\pi,\pi]^d$ 内では $t=0$ 以外にこの性質を持つ点はない．

逆に差の生成する部分群 $G$ が真部分群なら，
$\mathbb Z^d/G$ は非自明な有限生成Abel群である．
有限生成Abel群の構造定理により，
$\mathbb Z^d/G$ には非自明な $\mathbb T$ 値指標が存在する．
この指標は $t\in[-\pi,\pi]^d\setminus\{0\}$ による
$a\mapsto e^{it\cdot a}$ と表せ，$G$ 上で1である．
従って $e^{it\cdot X}$ は確率1で一定となり，
$|\varphi_X(t)|=1$ を得る．
\end{proof}

\begin{remark}
ここで「差の生成する群」が重要である．
値そのものが $\mathbb Z^d$ を生成するだけでは十分でない．
例えば $\{+1,-1\}$ の値を取る単純対称ランダムウォークでは
値の生成する群は $\mathbb Z$ だが，差の生成する群は
$2\mathbb Z$ である．
\end{remark}

\section{原点近傍の特性函数評価}
\begin{lemma}[二次展開]\label{lem:tay}
$X$ を $\mathbb R^d$ 値確率変数とし，
$\mathbb E|X|^2<\infty$ とする．平均を $\mu$，
共分散行列を $\Sigma$ とする．
$\psi(t)=e^{-it\cdot\mu}\mathbb E(e^{it\cdot X})$
とおけば
\[
 \psi(t)=1-\tfrac12t^T\Sigma t+o(|t|^2)
 \qquad(t\longrightarrow0).
\]
\end{lemma}
\begin{proof}
$Z=X-\mu$ とする．
\[
 e^{iu}=1+iu-\frac{u^2}{2}+u^2r(u)
\]
で $r(u)\to0$（$u\to0$）かつ $r$ が全実数上で有界となる
函数 $r$ を取れる．実際 $u=0$ では $r(0)=0$ とし，
$|u|\ge1$ で $|e^{iu}-1-iu+u^2/2|/u^2$ は有界である．
従って
\[
 \mathbb E(e^{it\cdot Z})
 =1-\tfrac12\mathbb E[(t\cdot Z)^2]
  +\mathbb E[(t\cdot Z)^2r(t\cdot Z)].
\]
$\mathbb E(t\cdot Z)=0$ であり
$\mathbb E[(t\cdot Z)^2]=t^T\Sigma t$ である．
最後の項を $|t|^2$ で割った絶対値は
$|Z|^2\sup|r|$ で支配され，
各 $Z$ に対して $r(t\cdot Z)\to0$ だから
優収束定理によって0に収束する．
\end{proof}

\begin{lemma}[小さい周波数でのGaussian減衰]\label{lem:bound}
$\Sigma$ が正定値なら，ある $\delta>0$ と $c>0$ が存在して
\[
 |\psi(t)|\le\exp(-c|t|^2)\qquad(|t|\le\delta)
\]
が成立する．
\end{lemma}
\begin{proof}
$\Sigma$ の最小固有値を $\lambda>0$ とする．
前補題により
$|\psi(t)|^2=1-t^T\Sigma t+o(|t|^2)$ である．
$\delta$ を十分小さく取れば
$|\psi(t)|^2\le1-\lambda|t|^2/2
\le\exp(-\lambda|t|^2/2)$ なので
$c=\lambda/4$ とすればよい．
\end{proof}

\begin{lemma}[GaussianのFourier積分]\label{lem:gauss}
$y\in\mathbb R^d$ に対して
\[
 (2\pi)^{-d}\int_{\mathbb R^d}
     e^{-iu\cdot y}e^{-|u|^2/2}\,du
       =(2\pi)^{-d/2}e^{-|y|^2/2}.
\]
\end{lemma}
\begin{proof}
各変数の積に分解できるから $d=1$ だけ示せばよい．
$I(y)=\int_{\mathbb R}e^{-u^2/2}e^{-iuy}\,du$ とする．
微分と積分の交換は $|u|e^{-u^2/2}$ の可積分性から従う．
部分積分を使うと
\[
 I'(y)=-i\int ue^{-u^2/2}e^{-iuy}\,du=-yI(y)
\]
である．よって $I(y)=I(0)e^{-y^2/2}$．
Gaussian積分 $I(0)=\sqrt{2\pi}$ を用いればよい．
後者は積分の二乗を極座標で計算して得られる．
\end{proof}

\section{多次元格子局所中心極限定理の証明}
\begin{proof}[定理\ref{thm:llclt}の証明]
$\varphi(t)=\mathbb Ee^{it\cdot X_1}$，
$\psi(t)=e^{-it\cdot\mu}\varphi(t)$ と置く．
Fourier反転により
\[
 \Pr(S_n=k)=(2\pi)^{-d}\int_{[-\pi,\pi]^d}
  e^{-it\cdot(k-n\mu)}\psi(t)^n\,dt.
\]
$x_{n,k}=\Sigma^{-1/2}(k-n\mu)/\sqrt n$ とし，
$u=\sqrt n\,\Sigma^{1/2}t$ と変数変換すると
\[
 n^{d/2}\sqrt{\det\Sigma}\,\Pr(S_n=k)
 =(2\pi)^{-d}\int_{D_n}
 e^{-iu\cdot x_{n,k}}\psi(\Sigma^{-1/2}u/\sqrt n)^n\,du,
\]
ただし $D_n=\sqrt n\,\Sigma^{1/2}[-\pi,\pi]^d$ とする．

$u$ が固定なら補題\ref{lem:tay}により
\[
 \psi(\Sigma^{-1/2}u/\sqrt n)^n
 \longrightarrow e^{-|u|^2/2}.
\]
補題\ref{lem:bound}から，
$|\Sigma^{-1/2}u/\sqrt n|\le\delta$ の領域では
\[
 |\psi(\Sigma^{-1/2}u/\sqrt n)|^n\le
  e^{-c|\Sigma^{-1/2}u|^2}\le e^{-c'|u|^2}
\]
となる．ここで $c'>0$ は $\Sigma$ の固有値にのみ依存する．
これによりこの領域での差の絶対値を，
$L^1(\mathbb R^d)$ 上の優収束定理で評価できる．

他方 $t\in[-\pi,\pi]^d$ で $|t|\ge\delta$ の部分では
強非周期性と特性函数の連続性から
$\rho=\max_{\delta\le|t|,\ t\in[-\pi,\pi]^d}|\varphi(t)|<1$ とできる．
したがってこの領域に対応する $u$ の積分は
$O(n^{d/2}\rho^n)$ であり0に収束する．
以上から，$\mathbf1_{D_n}(u)\psi(\Sigma^{-1/2}u/\sqrt n)^n$
は $L^1(\mathbb R^d)$ で $e^{-|u|^2/2}$ に収束する．
ここで複素値の優収束定理を実部・虚部に適用してよい．

$|e^{-iu\cdot x_{n,k}}|=1$ であるから，
両者のFourier積分の差の絶対値は
$L^1$ ノルムで一様に評価できる．
補題\ref{lem:gauss}を用いて
\[
 \sup_{k\in\mathbb Z^d}
 \left|n^{d/2}\sqrt{\det\Sigma}\,\Pr(S_n=k)
        -(2\pi)^{-d/2}e^{-|x_{n,k}|^2/2}\right|\to0
\]
を得る．これは定理の主張である．
\end{proof}

\section{一次元の一般格子と非格子密度}
\begin{corollary}[任意のspanをもつ整数分布]\label{cor:span}
$X_1,X_2,\ldots$ を独立同分布の整数値確率変数とし，
$0<\sigma^2=\operatorname{Var}(X_1)<\infty$ とする．
その台の差が生成する部分群を $q\mathbb Z$ とする
（$q$ は正整数）．$a\in\operatorname{supp}(X_1)$ を取る．
このとき $S_n$ は $na+q\mathbb Z$ 上に値を取り，
\[
 \sup_{k\in na+q\mathbb Z}\left|
   \frac{\sigma\sqrt n}{q}\,\Pr(S_n=k)
   -(2\pi)^{-1/2}
    \exp\!\left(-\frac{(k-n\mu)^2}{2n\sigma^2}\right)
 \right|\longrightarrow0.
\]
$k\notin na+q\mathbb Z$ では $\Pr(S_n=k)=0$ である．
\end{corollary}
\begin{proof}
$Y_i=(X_i-a)/q$ と置くと $Y_i\in\mathbb Z$ であり，
その台の差の最大公約数は1である．ゆえに定理\ref{thm:llclt}を
$Y_i$ に $d=1$ として適用できる．
$\operatorname{Var}(Y_i)=\sigma^2/q^2$，
$S_n=na+q\sum_iY_i$ を代入すればよい．
\end{proof}

\begin{lemma}[可積分な特性函数の無限遠での減衰]\label{lem:decay}
$X$ を $\mathbb R$ 値確率変数とし，
その特性函数 $\varphi$ が $L^1(\mathbb R)$ に属するとする．
このとき $\varphi(t)\to0$（$|t|\to\infty$）であり，
$X$ は有界連続密度
\[
 f(x)=\frac1{2\pi}\int_{\mathbb R}e^{-itx}\varphi(t)\,dt
\]
をもつ．
\end{lemma}
\begin{proof}
$\varphi$ は一様連続である．実際
$|\varphi(t+h)-\varphi(t)|\le\mathbb E|e^{ihX}-1|$ であり，
右辺は優収束定理により $h\to0$ で0へ収束する．
$\varphi$ が無限遠で0に収束しないなら，
$|\varphi(t_j)|\ge\varepsilon>0$，$|t_j|\to\infty$ とできる．
一様連続性によって一定幅の互いに素な区間上で
$|\varphi|\ge\varepsilon/2$ となり，$L^1$ 可積分性に矛盾する．

表示された $f$ は有界連続である．Fourier反転による
$\mathbb E g(X)=\int f(x)g(x)\,dx$ は
任意の $g\in C_c^\infty(\mathbb R)$ に対して
$g$ のFourier反転式とFubiniを用いて得られる．
$\int|\widehat g|<\infty$ だから積分交換は正当である．
測度の一意性から $f(x)dx$ が $X$ の分布を与え，
$f\ge0$ は連続性から，$\int f=1$ は確率測度の総質量から従う．
\end{proof}

\begin{theorem}[非格子密度の局所中心極限定理]\label{thm:density}
$X_i$ を独立同分布の実確率変数とし，
$0<\sigma^2=\operatorname{Var}(X_1)<\infty$，
$\mu=\mathbb EX_1$ とする．
$X_1$ は非格子型，すなわち
ある $a\in\mathbb R$，$h>0$ について
$\Pr(X_1\in a+h\mathbb Z)=1$ となることがないと仮定する．
さらに $X_1$ の特性函数 $\varphi$ が $L^1(\mathbb R)$ に
属すると仮定する．$S_n$ の連続密度を $f_n$ とすると
\[
 \sup_{x\in\mathbb R}\left|
 \sigma\sqrt n\,f_n(x)
 -(2\pi)^{-1/2}\exp\!\left(-\frac{(x-n\mu)^2}{2n\sigma^2}\right)
 \right|\longrightarrow0.
\]
\end{theorem}
\begin{proof}
前補題により各 $S_n$ は有界連続密度をもつ．
非格子性と補題\ref{lem:unit}から
$|\varphi(t)|<1$（$t\ne0$）である．
原点近傍では補題\ref{lem:tay}，\ref{lem:bound}を用いる．
$|t|\ge\delta>0$ 上では
$\varphi(t)\to0$（$|t|\to\infty$）と連続性により
$\rho=\sup_{|t|\ge\delta}|\varphi(t)|<1$ となる．
Fourier反転で $u=\sigma\sqrt n\,t$ と変数変換すると
\[
 \sigma\sqrt n\,f_n(x)
 =\frac1{2\pi}\int_{\mathbb R}
     e^{-iu(x-n\mu)/(\sigma\sqrt n)}
     \psi(u/(\sigma\sqrt n))^n\,du.
\]
原点近傍の項については前定理と同じGaussian減衰を用いる．
外側の積分は
\[
 \int_{|u|>\delta\sigma\sqrt n}
  |\psi(u/(\sigma\sqrt n))|^n\,du
 \le\sigma\sqrt n\,\rho^{n-1}
                      \int_{\mathbb R}|\varphi(t)|\,dt\to0
\]
である．従って積分核は
$e^{-u^2/2}$ に $L^1(\mathbb R)$ 収束する．
補題\ref{lem:gauss}とFourier積分の一様評価から結論を得る．
\end{proof}

\section{例と適用上の注意}
\begin{example}[二項分布]
$X_i\in\{0,1\}$ で $\Pr(X_i=1)=p$，$0<p<1$ とする．
台の差は $\mathbb Z$ を生成し，$\mu=p$，$\sigma^2=p(1-p)$ である．
従って $S_n\sim\operatorname{Bin}(n,p)$ に対し
\[
 \sup_{k\in\mathbb Z}\left|
 \sqrt{np(1-p)}\binom nk p^k(1-p)^{n-k}
 -(2\pi)^{-1/2}e^{-(k-np)^2/(2np(1-p))}
 \right|\to0.
\]
ここで $k\notin\{0,\ldots,n\}$ では二項係数項を0と解釈する．
\end{example}

\begin{example}[単純対称ランダムウォーク]
$X_i\in\{-1,+1\}$ で両者の確率が1/2とする．
台の差が生成する群は $2\mathbb Z$ であるから，
偶奇条件 $k\equiv n\pmod2$ の点に限り
\[
 \sup_{k\equiv n\ (\mathrm{mod}\,2)}
 \left|\frac{\sqrt n}{2}\Pr(S_n=k)
 -(2\pi)^{-1/2}e^{-k^2/(2n)}\right|\to0.
\]
偶奇条件を落とすと $k$ に一様な誤差評価は成立しない．
\end{example}

\begin{example}[多次元遅延ランダムウォーク]
$\mathbb Z^d$ 上で $X_i=0$ を確率1/2で取り，
$\pm e_j$（$j=1,\ldots,d$）をそれぞれ確率 $1/(4d)$ で取る．
この分布は強非周期的であり
$\mu=0$，$\Sigma=(2d)^{-1}I_d$ である．
従って定理\ref{thm:llclt}により特に
\[
 \Pr(S_n=0)\sim\left(\frac d{\pi n}\right)^{d/2}.
\]
\end{example}

\begin{remark}
一般格子分布では台の差の生成する部分群を無視できない．
非格子分布については本稿で $\varphi\in L^1$ という
十分条件を仮定した．この条件を仮定しない最も一般的な
非格子局所極限定理（区間確率に関するStone型の定理）は
本稿の証明の射程外である．また重尾分布の安定則や
多変量非格子型の一般定理は扱わない．
\end{remark}

\begin{thebibliography}{9}
\bibitem{feller} W. Feller, \emph{An Introduction to Probability Theory
 and Its Applications}, Vol. II, 2nd ed., Wiley, 1971.
\bibitem{petrov} V. V. Petrov, \emph{Sums of Independent Random Variables},
 Springer, 1975.
\bibitem{durrett} R. Durrett, \emph{Probability: Theory and Examples},
 5th ed., Cambridge University Press, 2019.
\end{thebibliography}

\section{到達課題の完全解答}
\noindent\textbf{対称ランダムウォークの原点確率．}
\begin{proposition}
$X_i$ を独立で $P(X_i=\pm1)=1/2$ とし，
$S_n=\sum_{i=1}^nX_i$ とする．
\[
P(S_{2m}=0)=\frac{\binom{2m}m}{2^{2m}}
\sim\frac1{\sqrt{\pi m}}
\]
が成り立つ．
\end{proposition}
\begin{proof}
$S_{2m}=0$ となるのは
$+1$ がちょうど $m$ 回現れる場合であり，
その選び方は $\binom{2m}m$ 通りである．
各列の確率は $2^{-2m}$ なので第一式が従う．
Stirlingの公式
$n!\sim\sqrt{2\pi n}(n/e)^n$ を用いれば
\[
\binom{2m}m
=\frac{(2m)!}{(m!)^2}
\sim\frac{4^m}{\sqrt{\pi m}}.
\]
これを $4^m$ で割れば結論を得る．
\end{proof}
\noindent\textbf{到達確認．}
$P(S_{2m+1}=0)$ を求めよ．
\par\noindent\textbf{解答．}
奇数個の $\pm1$ の和は奇数なので零にはならず，
確率は0である．格子の周期性に注意が必要である．
\end{document}